\par\Needspace{6\baselineskip}\medskip
\begin{center}\hypertarget{native-168-article121}{}\pdfbookmark[2]{121.}{article-121}121.\end{center}
Quoties itaque incognitarum valores approximati iam in potestate sunt,
veri inde per methodum modo explicatam omni quae desideratur praecisione
\SourceLeaf{169}derivari possunt. Primo scilicet computabuntur valores ipsarum $X,Y$ istis
valoribus approximatis $a,b$ respondentes: qui nisi sponte iam evanescunt,
calculus duobus aliis valoribus ab illis parum diversis $a',b'$ repetetur, ac dein
tertio systemate $a'',b''$, nisi fortuito ex secundo $X$ et $Y$ evanuerunt. Tunc
per formulas art. praec. valores veri elicientur, quatenus suppositio, cui illae
formulae innituntur, a veritate haud sensibiliter discrepat. De qua re quo
melius iudicium ferri possit, calculus valorum ipsarum $X,Y$ cum illis valoribus
correctis repetetur: qui si aequationibus $X=0$, $Y=0$ nondum satisfieri
monstrat, certe valores multo minores ipsarum $X,Y$ inde prodibunt, quam
per tres priores hypotheses, adeoque elementa orbitae hinc resultantia longe
exactiora erunt, quam ea, quae primis hypothesibus respondent. Quibus si
acquiescere nolumus, consultissimum erit, omissa ea hypothesi quae maximas
differentias produxerat, duas reliquas cum quarta denuo iungere, atque sic ad
normam art. praec. quintum systema valorum ipsarum $x,y$ formare: eodemque
modo, ubi operae pretium videbitur, ad hypothesin sextam etc. progredi
licebit, donec aequationibus $X=0$, $Y=0$ tam exacte satisfactum fuerit,
quam tabulae logarithmicae et trigonometricae permittunt. Rarissime tamen
opus erit, ultra systema quartum progredi, nisi hypotheses primae nimis adhuc a veritate aberrantes suppositae fuerint.
